\subsection{以 B × I 为底空间的局部平凡纤维空间}

命题 2. --- 设 B 是一个仿紧拓扑空间，且 (E, p) 是 B × I 上的局部平凡纤维空间。置 E₁ = \(\overline{p}^{-1}(B \times \{1\})\)，并记 \(p_{1}: E_{1} \to B\) 为映射 \(pr_{1} \circ p | E_{1}\)。则 B × I-空间 (E, p) 与 (E₁ × I, \(p_{1} \times Id_{I}\)) 同构。

我们先证明两个引理。

引理 2. --- 设 \(\alpha\)、\(\beta\)、\(\gamma\) 是满足 \(\alpha \leqslant \beta \leqslant \gamma\) 的实数；设 B 是一个拓扑空间，且 \(p: \mathrm{E} \to \mathrm{B} \times [\alpha, \gamma]\) 是连续映射。置 \(\mathrm{B}_0 = \mathrm{B} \times [\alpha, \beta]\)、\(\mathrm{B}_1 = \mathrm{B} \times [\beta, \gamma]\)、\(\mathrm{E}_0 = p^{-1}(\mathrm{B}_0)\)、\(\mathrm{E}_1 = p^{-1}(\mathrm{B}_1)\)，并记 \(p_0: \mathrm{E}_0 \to \mathrm{B}_0\)、\(p_1: \mathrm{E}_1 \to \mathrm{B}_1\) 为由 \(p\) 诱导的映射。若 \((\mathrm{E}_0, p_0)\) 和 \((\mathrm{E}_1, p_1)\) 是可平凡化的纤维空间，则 \((\mathrm{E}, p)\) 也具有同样性质。

令 \(g_0 \colon \mathrm{E}_0 \to \mathrm{B}_0 \times \mathrm{F}_0\) 和 \(g_1 \colon \mathrm{E}_1 \to \mathrm{B}_1 \times \mathrm{F}_1\) 分别为 \(\mathrm{E}_0\) 和 \(\mathrm{E}_1\) 的平凡化。记 \(g_0'\) 和 \(g_1'\) 为 \(\mathrm{B} \times \{\beta\}\)-纤维空间 \(p^{-1}(B \times \{\beta\})\) 的平凡化，它们由 \(g_0\) 和 \(g_1\) 通过限制导出。映射 \(h = g_0' \circ (g_1')^{-1}\) 是一个 \(B \times \{\beta\}\)-同构，将 \(B \times \{\beta\} \times F_1\) 映到 \(B \times \{\beta\} \times F_0\) 上。我们定义一个连续映射 \(h'\)，它从 \(B \times F_1\) 到 \(F_0\)，规定 \(h'(a, y) = \mathrm{pr}_3 \circ h(a, \beta, y)\) 对于 \((a, y) \in B \times F_1\)。对于 \((a, t, y) \in B \times [\beta, \gamma] \times F_1\)，置 \(H(a, t, y) = (a, t, h'(a, y))\)。这样定义的映射 H 是一个 \((B \times [\beta, \gamma])\)-同构，将 \(B \times [\beta, \gamma] \times F_1\) 映到 \(B \times [\beta, \gamma] \times F_0\) 上，并且有 \(g_0 | p^{-1}(B \times \{\beta\}) = H \circ g_1 | p^{-1}(B \times \{\beta\})\)。因而存在连续映射 \(g \colon E \to B \times [\alpha, \gamma] \times F_0\)，使得 \(g | E_0 = g_0\) 且 \(g | E_1 = H \circ g_1\)。映射 \(g\) 是一个 \(B \times [\alpha, \gamma]\)-空间同构，因此 E 是可平凡化的。

引理 3. --- 设 B 是一个拓扑空间，且 (E, p) 是 B×I 上的局部平凡纤维空间。B 的每一点 a 都有一个邻域 V，使得 V × I-空间 \(E_{V \times I}\) 是可平凡化的。

设 \(a\) 是 B 中的一点；对于每一点 \(t\) \(\mathbf{I}\)，存在 \(\mathrm{W}_t\)，它是 \(t\) 在 \(\mathbf{I}\) 中的一个开邻域，以及一个邻域 \(\mathrm{V}_t\)，它是 \(a\) 在 B 中的一个邻域，使得 E 在 \(\mathrm{V}_t \times \mathrm{W}_t\) 上可平凡化。因而存在一个整数 \(n > 0\)，并且对于每个整数 \(i\)（满足 \(1 \leqslant i \leqslant n\)），存在一点 \(t_i\)，它属于 \(\mathbf{I}\)，使得区间 \([\frac{i-1}{n}, \frac{i}{n}]\) 包含于 \(\mathrm{W}_{t_i}\) 中（III, p. 272, lemma 4）。置 \(\mathrm{V} = \cap_{1 \leqslant i \leqslant n} \mathrm{V}_{t_i}\)。纤维空间 E 在 \(\mathrm{V} \times [\frac{i-1}{n}, \frac{i}{n}]\) 上可平凡化，对于每个整数 \(i\)（满足 \(1 \leqslant i \leqslant n\)）。引理 3 随即由引理 2 对 \(n\) 作归纳得到。

现在证明命题。由引理 3，存在 B 的一个开覆盖 \((\mathrm{U}_j)_{j\in \mathrm{J}}\)，使得对于每个 \(j\in \mathrm{J}\)，E 在 \(\mathrm{U}_j\times \mathbf{I}\) 上可平凡化。由于空间 B 是仿紧的，我们可以设覆盖 \((\mathrm{U}_j)_{j\in \mathrm{J}}\) 是局部有限的（TG, IX, p. 49），并选取 B 的一个覆盖 \((\mathrm{A}_j)_{j\in \mathrm{J}}\)，使得对于每个 \(j\in \mathrm{J}\)，集合 \(\mathrm{A}_j\) 在 B 中是闭的且包含于 \(\mathrm{U}_j\) 中（TG, IX, p. 49, prop. 4 and p. 48, cor. 1）。

对于 B 的每个开子集 U，记 \(\mathcal{F}(\mathrm{U})\) 为所有 \(\mathrm{U}\times\mathbf{I}\) -同构的集合：这些同构将 \(\stackrel{-1}{p}(\mathrm{U}\times\mathbf{I})\) 映到 \(\stackrel{-1}{p_{1}}(\mathrm{U})\times\mathbf{I}\) 上，并将 \(\stackrel{-1}{p}(\mathrm{U}\times\{1\})\) 到 \(\stackrel{-1}{p_{1}}(\mathrm{U})\times\{1\}\) 的映射诱导为恒等映射。对于 B 的每一对开集 (V, U)

且满足 U ⊂ V，记 \(r_{\mathrm{UV}} \colon \mathcal{F}(\mathrm{V}) \to \mathcal{F}(\mathrm{U})\) 为如下映射：它将一个 \((\mathrm{V} \times \mathbf{I})\) -同构 \(g \colon \stackrel{-1}{p}(\mathrm{V} \times \mathbf{I}) \to \stackrel{-1}{p_0}(\mathrm{V}) \times \mathbf{I}\) 映到由 g 通过取子空间导出的 \((\mathrm{U} \times \mathbf{I})\) -同构。二元组 \(\mathcal{F} = ((\mathcal{F}(\mathrm{U})), (r_{\mathrm{UV}}))\) 是 B 上的一个层（I, p. 45, example 4）。为了证明命题，只需证明层 F 是柔软的（I, p. 64）。

取 \(j \in \mathrm{J}\)，令 A 为 \(\mathrm{A}_j\) 的闭子集，令 V 为 B 中的开集，满足 \(\mathrm{A} \subset \mathrm{V} \subset \mathrm{U}_j\)，并令 \(g\) 为 \(\mathcal{F}(\mathrm{V})\) 的一个元素。由于仿紧空间是正规的（TG, IX, p. 49, prop. 4），存在 A 的一个开邻域 W，使得 \(\overline{\mathrm{W}} \subset \mathrm{V}\)。我们将证明存在一个元素 \(g'\)，它属于 \(\mathcal{F}(\mathrm{U}_j)\)，使得 \(g' | \mathrm{W} = g | \mathrm{W}\)。于是，命题 6 的推论 2（I, p. 65）意味着层 \(\mathcal{F}\) 是柔软的。

由于 B × I-空间 E 在 \(U_{j} \times I\) 上可平凡化，我们可以设 \(\overline{p}^{-1}(U_{j} \times I) = U_{j} \times I \times F\)，其中 F 是一个拓扑空间。于是，层 \(\mathcal{F}(V)\) 的元素 g 是 V × I 到自身的 V × I-同构，它在 V × \{1\}×F 上诱导恒等映射。将 cor. 5 of IV, p. 439 应用于空间 \(X' = X\)、\(X = U_{j}\)、A = overline{W}、\(U = V\) 和 \(Y = Z = F\)，应用于映射 g: V 	imes I 	imes F 	o V 	imes I 	imes F，以及 \(U_{j} \times \{1\} \times F\) 的恒等映射。因而存在一个 \(U_{j} \times I\)-同构 \(g'\)，它将 \(U_{j} \times I \times F\) 映到自身，在 \(U_{j} \times \{1\} \times F\) 上诱导恒等映射，并且在 \(\overline{W} \times I \times F\) 上与 g 重合，从而在 W × I × F 上也重合。因此命题得证。

推论 1. --- 设 B 是一个仿紧拓扑空间，且 (E, p) 是 B × I 上的局部平凡纤维空间（I, p. 71, corollary 2）。对于 \(t \in I\)，记 \((E_{t}, p_{t})\) 为局部平凡的纤维 B-空间 \(i_{t}^{*}E\)，其中 \(i_{t} \colon B \to B \times \{t\}\) 是由 \(b \mapsto (b, t)\) 给出的映射。局部平凡的纤维 B-空间 \(E_{0}\) 和 \(E_{t}\) 对每个 \(t \in I\) 都同构。

推论 2. --- 设 B 和 B' 是拓扑空间，E 是 B 上的局部平凡纤维丛。令 \(f_0\) 和 \(f_1\) 为从 B' 到 B 的连续映射。设空间 B' 是仿紧的。若映射 \(f_0\) 和 \(f_1\) 同伦，则局部平凡的纤维 B'-空间 \(f_0^*E\) 和 \(f_1^*E\) 同构。

令 \(\sigma\colon\mathrm{B}^{\prime}\times\mathbf{I}\to\mathrm{B}\) 为连接 \(f_{0}\) 与 \(f_{1}\) 的同伦，并记 \(\mathrm{E}^{\prime}\) 为 \(\mathrm{B}^{\prime}\times\mathbf{I}\)-空间 \(\sigma^{*}\mathrm{E}\)；它是一个局部平凡纤维空间。记 \(i_{0}\) 和 \(i_{1}\) 为从 \(\mathrm{B}^{\prime}\) 到 \(\mathrm{B}^{\prime}\times\mathbf{I}\) 的映射，它们分别由 \(x\mapsto(x,0)\) 和 \(x\mapsto(x,1)\) 给出。由 Cor. 1，\(\mathrm{B}^{\prime}\)-纤维空间 \(i_{0}^{*}\mathrm{E}^{\prime}\) 和 \(i_{1}^{*}\mathrm{E}^{\prime}\) 同构。由于 \(\sigma\circ i_{0}=f_{0}\)，\(\mathrm{B}^{\prime}\)-空间 \(i_{0}^{*}\mathrm{E}^{\prime}\) 被识别为 \(f_{0}^{*}\mathrm{E}\)；

同样，B'-空间 \(i_{1}^{*}E'\) 被识别为 \(f_{1}^{*}E\)。因此，B'-纤维空间 \(f_{0}^{*}E\) 和 \(f_{1}^{*}E\) 同构，这正是要证明的结论。

推论 3. --- 设 B 是一个仿紧拓扑空间。若 B 与一个点同伦等价，则每个以 B 为底的局部平凡纤维丛都是可平凡化的。

推论 4. --- 设 B 和 B' 是拓扑空间，(E, p) 是 B 上的局部平凡纤维空间，f 是从 B' 到 B 的连续映射，\(\sigma: \mathrm{B}' \times \mathbf{I} \to \mathrm{B}\) 是以 f 为初始映射的同伦，且 \(\widetilde{f}\) 是 f 到 E 的连续提升。若空间 B' 是仿紧的，则存在同伦 \(\widetilde{\sigma}: \mathrm{B}' \times \mathbf{I} \to \mathrm{E}\)，它以 \(\widetilde{f}\) 为初始映射，并且是 \(\sigma\) 到 E 的提升。

设 \((\mathrm{E}^{\prime},p^{\prime})\) 为从 \(B^{\prime}\times I\)-空间 \((\mathrm{E},p)\) 沿映射 \(\sigma\colon B^{\prime}\times I\to B\) 进行基变换得到的空间。它是一个局部平凡纤维空间（I, p. 71, cor. 2），且映射 \(s\colon B^{\prime}\times\{0\}\to E^{\prime}\) 由 \(s((a,0))=((a,0),\widetilde{f}(a))\) 定义；对于 \(a\in B^{\prime}\)，它是连续的，并且是 \(p^{\prime}\) 在 \(B^{\prime}\times\{0\}\) 上的连续截面。由 Prop. 2，存在一个连续截面 \(\widetilde{s}\)，它是 \(p^{\prime}\) 的截面，并延拓 s。映射 \(\widetilde{\sigma}=\operatorname{pr}_{2}\circ\widetilde{s}\) 具有所需性质。

